\section*{Limitations}
\label{sec:limitations}
% ACL requires an unnumbered Limitations section; it does not count toward the
% page limit.
% OWNER: Sanjana Garimella (pairs with discussion)

\paragraph{Synthetic, uniform noise.}
Our corruption is a fixed per-sentence probability of flipping the verb's number
feature, applied independently and uncorrelated with any property of the
sentence. Naturalistic agreement errors are not like this: they cluster by
speaker, register, and construction, and arise disproportionately in the hard
configurations (long dependencies and intervening attractors), where human
production itself slips. A model trained on realistically structured noise might
degrade along a different curve, perhaps eroding on attractor items sooner than
our uniform manipulation predicts. Our noise also targets only the verb's number;
we do not model omissions, wrong lemmas, or agreement with coordinated subjects,
so the study isolates one axis of ungrammaticality rather than spanning the
natural distribution of errors.

\paragraph{Templatic, vocabulary-controlled corpus.}
The training and evaluation data come from a controlled subject--verb agreement
corpus with a restricted vocabulary and lowercased, anonymized tokens. This
control is what makes the leak-free minimal-pair design and the attractor-gap
metric clean, but it bounds external validity: the sentences are simpler and more
template-like than running text, and the absolute accuracies, along with the
threshold near $50\%$ noise, may not transfer to models trained on naturalistic
corpora. We therefore read every comparison as a relative, within-setup contrast
between architectures, not as a calibrated estimate of either model's behavior on
natural language.

\paragraph{RNNG ceiling is bounded by parse quality.}
The RNNG is trained on benepar parses of the clean corpus rather than gold trees.
On lowercased, anonymized, templatic text the parser makes systematic errors,
and those errors propagate into the structures the RNNG learns from. Its observed
robustness is thus a lower bound on what a correctly-parsed structural model could
achieve, and we cannot fully rule out that part of the RNNG--LSTM gap reflects
parser idiosyncrasies rather than structural bias per se. A model with a
hierarchical inductive bias but no dependence on an external parser (e.g.\ an
ON-LSTM) would be an informative control that we leave to future work.

\paragraph{Scale, seeds, and training length.}
The architecture comparison currently rests on a single full-scale RNNG seed.
That seed is scored on the full held-out set, which resolves the earlier coverage
mismatch with the LSTM, but we cannot yet attach error bars or a significance
test to the cross-architecture contrast, and the multi-seed replication (seeds~1
and~3) is still in progress. Both models are also modest in scale (two layers,
$150$k training sentences, five to eight epochs). We confirm that the LSTM
learning trajectories converge, but we do not exhaustively verify RNNG
convergence at eight epochs; larger models or longer training could shift the
breaking point and should temper any precise reading of the crossover noise rate.
Finally, the two architectures' \emph{overall} accuracies are compared on
near-identical rather than identical held-out sets (a $\sim$99\% overlap,
$19{,}983$ vs.\ $19{,}819$ pairs, owing to which long sentences each scorer
covers); the within-model attractor gap, a rate, is unaffected, and an exact
intersection-matched comparison is left to future work.

\paragraph{Single language and phenomenon.}
All results concern English subject--verb number agreement. English has
impoverished agreement morphology and rigid word order, which makes the linear
nearest-noun heuristic an unusually strong competitor to the hierarchical rule.
Languages with richer agreement or freer word order, or other long-distance
dependencies such as reflexive binding or filler--gap constructions, could yield
a different noise-robustness profile and even a different ordering of the two
architectures. Whether ``structure buys robustness'' generalizes beyond English
subject--verb agreement remains an open question.
